\section{Leakage kernels: decay and evaporation as exit from tick traps (extended interface module)}
\label{app:leakage_kernel}

\paragraph{Scope and status.}
\AuditTag This appendix defines a finite-family ``leakage kernel'' vocabulary that unifies decay and evaporation as exits from tick traps.
It is an interface/audit module: it introduces no new theorem-level premise into the folding core.
Its role is to provide reproducible, CAP-auditable diagnostics (delay proxies, lifetimes, exit-channel decompositions) used by the relative-horizon language (Appendix~\ref{app:protocol_horizon_tick_trap}) and by the low-leakage phase module (Appendix~\ref{app:protected_low_leakage_phase}).

\subsection{Survival kernels and leakage rates}
\label{subsec:leakage_survival_kernels}

\paragraph{Finite-family modeling discipline.}
\AuditTag A leakage kernel is specified by choosing a finite candidate family of survival laws and a deterministic tie-break, then selecting the CAP-minimal element relative to declared mismatch costs.
This ensures that any numerical rows shown in generated tables are reproducible and do not hide untracked degrees of freedom.

\begin{definition}[Survival kernel (finite-family)]
\label{def:survival_kernel_family}
Fix a finite index set $\mathcal F$ and a family of survival functions $\{P_f(t)\}_{f\in\mathcal F}$ with $P_f(0)=1$ and $P_f(t)\in[0,1]$.
The associated leakage rate proxy is defined by
\[
\Gamma_f := \inf\{\gamma>0:\ P_f(t)\le \e^{-\gamma t}\ \text{for all }t\ge 0\},
\]
and a lifetime proxy by $\tau_f := 1/\Gamma_f$ when $\Gamma_f>0$.
\end{definition}

\begin{remark}[Exponential law as a reference]
\label{rem:exponential_reference}
If $P(t)=\e^{-\Gamma t}$, then $\Gamma$ is the standard decay width and $\tau=1/\Gamma$ is the mean lifetime.
This appendix does not assume an exponential law universally; it only uses a finite candidate family.
\end{remark}

\subsection{A minimal finite-family demo (reproducible audit)}
\label{subsec:leakage_kernel_demo}

\AuditTag The deterministic generator \path{scripts/exp\_leakage\_kernel\_demo.py} provides a compact, finite-family illustration of the survival-kernel vocabulary (candidate family, CAP selection, and reported proxies).

\begin{table}[t]
\centering
\scriptsize
\setlength{\tabcolsep}{6pt}
\renewcommand{\arraystretch}{1.15}
\begin{tabular}{l r r l}
\toprule
family id & $\Gamma$ (proxy) & $\tau$ (proxy) & note \\
\midrule
exp\_fast & 1.000000 & 1.000000 & exponential; high leakage \\
exp\_mid & 0.100000 & 10.000000 & exponential; mid leakage \\
exp\_slow & 0.010000 & 100.000000 & exponential; low leakage; CAP-selected \\%
\\
\bottomrule
\end{tabular}
\caption{A minimal leakage-kernel demo table reproduced by \texttt{scripts/exp\_leakage\_kernel\_demo.py}.}
\label{tab:leakage_kernel_demo}
\end{table}

\subsection{Delay as the trap readout}
\label{subsec:leakage_delay_readout}

\paragraph{Delay proxy.}
\InterfaceTag As recorded in Section~\ref{sec:mass_latency_coordinate}, an operational proxy for protocol overhead is Wigner--Smith time delay.
We treat a large delay proxy as an indicator of a deep trap (large stabilizing overhead), and a small leakage rate as an indicator of a well-protected trap.
The delay proxy is complementary to the lifetime proxy: both are reported as audit outputs.
\AuditTag When delay is extracted from a scattering readout, the algebraic bridge to the manuscript's determinant/trace vocabulary is the logdet identity:
on a declared unitary band, $\Tr Q(\omega)=-\iu\,\dd(\log\det S(\omega))/\dd\omega$ (Lemma~\ref{lem:wigner_smith_trace_logdet}).
Accordingly, the minimal failure point is the same as for scattering-delay audits: approximate unitarity on the declared band (S1; Appendix~\ref{app:minimal_failure_point_templates}).

\subsection{Exit-channel decompositions}
\label{subsec:exit_channel_decomposition}

\begin{definition}[Exit channels (protocol)]
\label{def:exit_channels_protocol}
Fix a coarse-graining that exposes a finite set of externally distinguishable exit channels $\mathcal C$.
An exit-channel decomposition for a target $X$ is a probability vector
\[
\mathbf p(X)=(p_c(X))_{c\in\mathcal C},\qquad p_c(X)\ge 0,\ \sum_{c\in\mathcal C}p_c(X)=1,
\]
representing the relative weights of leakage into each channel under the chosen coarse-graining.
\end{definition}

\paragraph{Interpretation.}
\AuditTag In this vocabulary, ``radiation'' is not a primitive label; it is leakage into some externally propagating subset of channels.
Which subset is interpreted as ``radiative'' depends on additional matching assumptions (e.g.\ symmetry breaking, propagation conditions), and is therefore not fixed by protocol structure alone.

\subsection{The \texorpdfstring{$m=6$}{m=6} specialization: \texorpdfstring{$18$}{18} trap states and \texorpdfstring{$3$}{3} exit channels}
\label{subsec:m6_trap_exit}

\paragraph{Trap/exit split at \texorpdfstring{$m=6$}{m=6}.}
\AuditTag At window length $m=6$ the stable-type decomposition has $18$ cyclic types and $3$ boundary types.
In the present appendix we treat the $18$ cyclic types as \emph{trap states} (discrete trap categories) and the $3$ boundary types as a distinguished \emph{exit-channel family}.
This is an audit choice: it does not identify a unique microscopic mechanism; it identifies a clean, finite protocol-level channelization.

\paragraph{Exit channels as gauge-sector channels (not ``photons'').}
\AuditTag In the SM labeling closure, the three boundary types at $m=6$ are uniquely and monotonically assigned to $U(1)\prec SU(2)\prec SU(3)$ (Remark~\ref{rem:boundary_value_order}).
Therefore, in this specialization, the $3$ exit channels should be read as a \emph{gauge-sector channel family} rather than as a literal identification with a single radiative species.
Interpreting a particular channel as ``electromagnetic radiation'' requires additional matching assumptions beyond the protocol structure.

\paragraph{Reproducible audit rows.}
\AuditTag The following table rows (generated by a deterministic script) report, for each of the $18$ trap categories, a declared delay proxy level and an exit-channel weight vector into the three boundary channels.
\begin{table}[t]
\centering
\scriptsize
\setlength{\tabcolsep}{6pt}
\renewcommand{\arraystretch}{1.15}
\begin{tabular}{l l r r r}
\toprule
trap type $w$ & label & delay level & $p_{U(1)}$ & $(p_{SU(2)},p_{SU(3)})$ \\
\midrule
000000 & $\nu_R^{(1)}$ & 1 & 0.387 & (0.323,0.290) \\
100000 & $L_L^{(1)}$ & 1 & 0.391 & (0.322,0.288) \\
010000 & $e_R^{(1)}$ & 1 & 0.395 & (0.321,0.285) \\
001000 & $Q_L^{(1)}$ & 1 & 0.399 & (0.319,0.282) \\
101000 & $d_R^{(1)}$ & 1 & 0.404 & (0.318,0.278) \\
000100 & $u_R^{(1)}$ & 1 & 0.411 & (0.316,0.274) \\
100100 & $\nu_R^{(2)}$ & 2 & 0.418 & (0.313,0.269) \\
010100 & $e_R^{(2)}$ & 2 & 0.427 & (0.310,0.263) \\
000010 & $u_R^{(2)}$ & 2 & 0.438 & (0.307,0.255) \\
100010 & $L_L^{(2)}$ & 2 & 0.452 & (0.301,0.247) \\
010010 & $Q_L^{(2)}$ & 2 & 0.471 & (0.294,0.235) \\
001010 & $\nu_R^{(3)}$ & 2 & 0.496 & (0.283,0.220) \\
101010 & $L_L^{(3)}$ & 3 & 0.533 & (0.267,0.200) \\
000001 & $d_R^{(2)}$ & 3 & 0.593 & (0.237,0.169) \\
010001 & $e_R^{(3)}$ & 3 & 0.484 & (0.323,0.194) \\
001001 & $Q_L^{(3)}$ & 3 & 0.308 & (0.462,0.231) \\
000101 & $d_R^{(3)}$ & 3 & 0.167 & (0.417,0.417) \\
010101 & $u_R^{(3)}$ & 3 & 0.160 & (0.280,0.560) \\%
\\
\bottomrule
\end{tabular}
\caption{Audit rows for the $m=6$ trap/exit specialization: $18$ cyclic trap categories and $3$ boundary exit channels (assigned to $U(1),SU(2),SU(3)$ by the SM labeling closure). Rows are reproduced by \texttt{scripts/exp\_leakage\_kernel\_m6\_trap\_exit.py}.}
\label{tab:leakage_kernel_m6_trap_exit}
\end{table}

\paragraph{Audit summary (m=6 trap/exit table).} \AuditTag This fragment treats the 18 cyclic stable types at $m=6$ as trap categories and the 3 boundary types as exit channels. Exit weights are computed by intrinsic-value proximity to the boundary intrinsic values $V(100001)=14$, $V(101001)=17$, $V(100101)=19$ (assigned to $U(1),SU(2),SU(3)$ by the SM labeling closure). A finite family over exponent $p\in\{1,2,3\}$ is CAP-selected by minimizing a monotonicity-violation cost; the selected value is $p^\ast=1$ with cost 11.235314.%

